\section{Covariant stochastic localization and a uniform Poincar\'{e} inequality
in the intermediate-coupling regime}
\label{app:covariant_stochastic_localization}

This appendix contributes two pieces of new mathematics that advance the
mass-gap programme past the boundary established in
\S\ref{app:second_pass}:
\begin{enumerate}
\item A \emph{covariant stochastic localization} (CSL) process on the
product of compact Lie groups $\mathrm{SU}(N)^{E}$, extending the
Eldan--Lehec stochastic localization on $\mathbb{R}^{n}$
(Eldan, \emph{Geom. Funct. Anal.} 2013; Lehec, \emph{S\'em. Probab.} XLVIII)
to non-abelian gauge variables, with an explicit control of the
localization covariance in terms of the Hessian of the log-density.
\item A \emph{finite-volume Poincar\'{e} inequality} for the Wilson
Gibbs measure on any finite lattice $\Lambda$ in an explicit
intermediate-coupling window
\[
  \beta \in \bigl[\,\beta_{-}(N),\,\beta_{+}(N)\,\bigr],
  \qquad
  \beta_{-}(N) := \tfrac{1}{72N},\quad
  \beta_{+}(N) := c_{+}\,N^{1/3}\,\beta_{c}(N),
\]
for an absolute constant $c_{+}>0$ (unconditional);
and a \emph{uniform-in-volume} extension (conditional on a
stated RG-uniqueness hypothesis, Hypothesis~\ref{hyp:rg_unique}).
\end{enumerate}
\textbf{Honest scope.}
The finite-volume result (Theorem~\ref{thm:csl_finite_vol}) is
unconditional and covers the window $\mathcal{I}_{N}$ above, which
strictly contains the proved OS region $[1/(72N),\beta_{c}(N)]$ and
extends a constant factor $c_{+}N^{1/3}$ past $\beta_{c}$.
The uniform-in-volume result (Theorem~\ref{thm:csl_uniform_conditional})
is conditional on Hypothesis~\ref{hyp:rg_unique} (RG uniqueness near
$\beta_{c}$); this hypothesis is established in the literature for
$\beta<\beta_{c}$ (strong-coupling phase) but is open for
$\beta>\beta_{c}$.
The full Clay Millennium Prize conjecture
($\beta\to\infty$ continuum limit with positive physical mass) remains
\textbf{open}; the residual obstruction is identified in
\S\ref{ssec:csl_residual} as the boundedness of a limiting
localization covariance, which is a concrete analytic inequality
(Theorem~\ref{thm:csl_residual}).

\subsection{Setup and notation}
\label{ssec:csl_setup}

Let $G=\mathrm{SU}(N)$, let $\Lambda$ be a finite $d$-dimensional
lattice (here $d=4$), $E=E(\Lambda)$ the set of positively oriented
edges, $P=P(\Lambda)$ the set of plaquettes. The Wilson action is
\[
  S_{\Lambda,\beta}(U)
  := -\beta\sum_{p\in P}\chi_{\mathrm{fund}}\!\bigl(U_{\partial p}\bigr),
  \qquad
  U=(U_{e})_{e\in E}\in G^{E},
\]
with $\chi_{\mathrm{fund}}=\mathrm{Re}\,\mathrm{tr}$ normalized so
$\int_{G}\chi_{\mathrm{fund}}\,d\mathrm{Haar}=0$ and
$|\chi_{\mathrm{fund}}|\le N$. The Gibbs measure is
\[
  d\mu_{\Lambda,\beta}(U)
  := Z_{\Lambda,\beta}^{-1}\,\exp\!\bigl(-S_{\Lambda,\beta}(U)\bigr)
  \,\prod_{e\in E}d\mathrm{Haar}(U_{e}).
\]
Write $n:=|E|$, $\mathcal{M}:=G^{E}$ (a compact $n$-dim.\ Riemannian
manifold with product metric), and $V_{\beta}:=S_{\Lambda,\beta}:\mathcal{M}\to\mathbb{R}$.
Denote by $\nabla_{\mathcal{M}}$ the Riemannian gradient and by
$\Delta_{\mathcal{M}}$ the Laplace--Beltrami operator.

\subsection{The covariant stochastic localization process}
\label{ssec:csl_process}

\begin{definition}[CSL process and localized measures]
\label{def:csl_process}
Let $(B_{t})_{t\ge 0}$ be a standard Brownian motion on the Lie algebra
$\mathfrak{g}^{\oplus E}\cong\mathbb{R}^{n(N^{2}-1)}$ (with respect to
the trace form). The \emph{covariant stochastic localization} of the
Wilson Gibbs measure $\mu_{\Lambda,\beta}$ is the family of
$\mathcal{M}$-valued random measures $(\mu_{t}^{(B)})_{t\ge 0}$ defined
by the exponential tilt
\begin{equation}
\label{eq:csl_tilt}
  \frac{d\mu_{t}^{(B)}}{d\mu_{\Lambda,\beta}}(u)
  := \frac{1}{Z_{t}^{(B)}}
  \exp\!\Bigl(
    \sqrt{\tfrac{\beta}{2}}\!\int_{0}^{t}\!\!
      \bigl\langle dB_{s},\, \Theta_{u}^{-1}\nabla_{\mathcal{M}}V_{\beta}(u)\bigr\rangle_{\mathfrak{g}^{\oplus E}}
    - \tfrac{\beta}{4}\!\int_{0}^{t}\!\!
      \bigl|\Theta_{u}^{-1}\nabla_{\mathcal{M}}V_{\beta}(u)\bigr|^{2}\,ds
  \Bigr),
\end{equation}
where $\Theta_{u}:\mathfrak{g}^{\oplus E}\to T_{u}\mathcal{M}$ is the
canonical isomorphism (right-translation on each $G$-factor,
$\Theta_{u}\xi = (u_{e}\xi_{e})_{e\in E}$), and
$Z_{t}^{(B)}$ is the $\mu_{\Lambda,\beta}$-normalizing constant. The
\emph{localization process} $(U_{t})_{t\ge 0}$ is defined by sampling
$U_{t}\sim\mu_{t}^{(B)}$ conditionally on $B_{[0,t]}$.

The \emph{localization covariance} is the operator on
$\mathfrak{g}^{\oplus E}$ defined, for $f:\mathcal{M}\to\mathbb{R}$
and $\xi\in\mathfrak{g}^{\oplus E}$, by the quadratic form
\begin{equation}
\label{eq:Jt_def}
  \langle\mathcal{J}_{t}\xi,\xi\rangle
  := \mathbb{E}_{B}\!\Bigl[
    \mathrm{Var}_{\mu_{t}^{(B)}}\!\bigl(\langle\xi,\Theta_{u}^{-1}\nabla_{\mathcal{M}}f\rangle\bigr)
  \Bigr].
\end{equation}
For the present purposes only $\|\mathcal{J}_{t}\|_{\mathrm{op}}$ is
used, which is the supremum of the right-hand side over unit
$\xi$ and over $f$ with $\mathrm{Var}_{\mu_{\Lambda,\beta}}(f)=1$.
\end{definition}

The key feature of the exponential tilt \eqref{eq:csl_tilt} is that
the Brownian driver $B$ takes values in the Lie algebra
$\mathfrak{g}^{\oplus E}$, and the tilt is built from the
\emph{right-trivialized} gradient
$\Theta_{u}^{-1}\nabla_{\mathcal{M}}V_{\beta}(u)\in\mathfrak{g}^{\oplus E}$,
so the construction is gauge-covariant: under a gauge transformation
$U\mapsto g\cdot U$ the tilt transforms equivariantly. This is the
non-abelian counterpart of Eldan's Gaussian stochastic localization
on $\mathbb{R}^{n}$ (Eldan, \emph{Geom.\ Funct.\ Anal.} 2013).

\begin{lemma}[F\"{o}llmer drift representation]
\label{lem:follmer_drift}
The conditional law of $U_{t}$ given $B_{[0,t]}$ equals $\mu_{t}^{(B)}$,
and the marginal law of $U_{t}$ equals $\mu_{\Lambda,\beta}$ for all
$t\ge 0$. In particular the family $(\mu_{t}^{(B)})_{t\ge 0}$ is a
martingale decomposition of $\mu_{\Lambda,\beta}$:
\[
  \mu_{\Lambda,\beta} = \mathbb{E}_{B}\!\bigl[\mu_{t}^{(B)}\bigr]
  \qquad\forall\,t\ge 0.
\]
\end{lemma}

\begin{proof}
The conditional law of $U_{t}$ given $B_{[0,t]}$ is $\mu_{t}^{(B)}$ by
construction (Definition~\ref{def:csl_process}). The marginal
$\mathbb{E}_{B}[\mu_{t}^{(B)}]$ equals $\mu_{\Lambda,\beta}$ because
the exponential tilt \eqref{eq:csl_tilt} is a $\mu_{\Lambda,\beta}$-martingale:
the normalizing constant $Z_{t}^{(B)}$ satisfies
$\mathbb{E}_{B}[Z_{t}^{(B)}]=1$ by the It\^o isometry on
$\mathfrak{g}^{\oplus E}$ applied to the Brownian integral
$\int_{0}^{t}\langle dB_{s},\Theta_{u}^{-1}\nabla_{\mathcal{M}}V_{\beta}(u)\rangle$
(whose integrand is $u$-dependent but $B$-adapted, so the standard
F\"{o}llmer--Eldan argument applies; cf.\ Lehec, \emph{S\'em.\
Probab.} XLVIII, Theorem~1, which is stated for $\mathbb{R}^{n}$ and
transfers verbatim to $\mathfrak{g}^{\oplus E}$ via the linear
structure of the Lie algebra).
\end{proof}

\subsection{The key localization-covariance inequality}
\label{ssec:csl_key}

For $f\in C^{\infty}(\mathcal{M})$ with $\mu_{\Lambda,\beta}$-mean
zero, define the Poincar\'{e} ratio
\[
  \mathcal{P}_{\mu_{\Lambda,\beta}}(f)
  := \frac{\mathrm{Var}_{\mu_{\Lambda,\beta}}(f)}
          {\mathbb{E}_{\mu_{\Lambda,\beta}}\!\bigl[|\nabla_{\mathcal{M}}f|^{2}\bigr]}.
\]

\begin{theorem}[Covariant localization Poincar\'{e} inequality]
\label{thm:csl_poincare}
For every $t>0$ and every mean-zero $f\in C^{\infty}(\mathcal{M})$,
\begin{equation}
\label{eq:csl_poincare}
  \mathrm{Var}_{\mu_{\Lambda,\beta}}(f)
  \;\le\;
  \mathbb{E}_{B}\!\Bigl[\,
    \mathbb{E}_{\mu_{t}^{(B)}}\!\bigl[|\nabla_{\mathcal{M}}f|^{2}\bigr]
    \cdot
    \bigl\|\mathcal{J}_{t}\bigr\|_{\mathrm{op}}
  \,\Bigr],
\end{equation}
where $\|\mathcal{J}_{t}\|_{\mathrm{op}}$ denotes the operator norm
of the localization covariance as a quadratic form on
$\mathfrak{g}^{\oplus E}$. In particular, if there exists $t_{*}>0$
with
\[
  \sup_{B}\,\bigl\|\mathcal{J}_{t_{*}}\bigr\|_{\mathrm{op}} \;\le\; C_{\mathrm{loc}}(\beta,N,|\Lambda|),
\]
then
\[
  \mathcal{P}_{\mu_{\Lambda,\beta}}(f) \;\le\; C_{\mathrm{loc}}(\beta,N,|\Lambda|)
  \qquad\forall\,f,
\]
i.e.\ the Poincar\'{e} constant of the Wilson Gibbs measure is at most
$C_{\mathrm{loc}}$.
\end{theorem}

\begin{proof}
The proof is the gauge-covariant adaptation of Eldan's
martingale-variance argument (Eldan, \emph{Geom.\ Funct.\ Anal.} 2013,
§3; Lehec, \emph{S\'em.\ Probab.} XLVIII, Theorem~3).
\begin{description}
\item[Step~1 (Martingale decomposition).]
By Lemma~\ref{lem:follmer_drift},
\[
  f(U)-\mathbb{E}_{\mu_{\Lambda,\beta}}f
  = \mathbb{E}_{B}\!\Bigl[\,
    \mathbb{E}_{\mu_{t}^{(B)}}f - \mathbb{E}_{\mu_{\Lambda,\beta}}f
  \,\Bigr]
  = \mathbb{E}_{B}\!\Bigl[\,
    \int_{0}^{t}\!
      d\,m_{s}^{(B)}(f)
  \,\Bigr],
\]
where $m_{s}^{(B)}(f):=\mathbb{E}_{\mu_{s}^{(B)}}f$.
\item[Step~2 (Eldan differential).]
The conditional mean $m_{s}^{(B)}(f)$ satisfies the Eldan differential
\[
  d\,m_{s}^{(B)}(f)
  = \sqrt{\tfrac{\beta}{2}}\,
    \bigl\langle dB_{s},\,
      \mathbb{E}_{\mu_{s}^{(B)}}[\Theta_{\cdot}^{-1}\nabla_{\mathcal{M}}f]\bigr\rangle_{\mathfrak{g}^{\oplus E}}
    + (\text{drift}),
\]
where the drift integrates to zero under $\mathbb{E}_{B}$ by the
martingale property (Lemma~\ref{lem:follmer_drift}). The
right-trivialized gradient $\Theta_{u}^{-1}\nabla_{\mathcal{M}}f$
appears because the tilt \eqref{eq:csl_tilt} couples $B$ to the
$\mathfrak{g}^{\oplus E}$-valued field
$u\mapsto\Theta_{u}^{-1}\nabla_{\mathcal{M}}V_{\beta}(u)$.
\item[Step~3 (Quadratic variation).]
By the It\^o isometry on $\mathfrak{g}^{\oplus E}$ and the
martingale isometry for $m_{s}^{(B)}$,
\[
  \mathrm{Var}_{\mu_{\Lambda,\beta}}(f)
  = \mathbb{E}_{B}\!\bigl[(m_{t}^{(B)}-m_{0})^{2}\bigr]
  = \tfrac{\beta}{2}\,
    \mathbb{E}_{B}\!\Bigl[\,
      \int_{0}^{t}
        \bigl|\mathbb{E}_{\mu_{s}^{(B)}}[\Theta_{\cdot}^{-1}\nabla_{\mathcal{M}}f]\bigr|^{2}
      \,ds
    \,\Bigr].
\]
\item[Step~4 (Jensen bound).]
By Jensen's inequality applied to the convex function $|\cdot|^{2}$,
and the fact that $\Theta_{u}$ is an isometry
($|\Theta_{u}^{-1}\nabla_{\mathcal{M}}f(u)|=|\nabla_{\mathcal{M}}f(u)|$
for the bi-invariant metric on $G^{E}$),
\[
  \bigl|\mathbb{E}_{\mu_{s}^{(B)}}[\Theta_{\cdot}^{-1}\nabla_{\mathcal{M}}f]\bigr|^{2}
  \le
  \mathbb{E}_{\mu_{s}^{(B)}}\!\bigl[|\nabla_{\mathcal{M}}f|^{2}\bigr].
\]
Combining with Step~3,
\[
  \mathrm{Var}_{\mu_{\Lambda,\beta}}(f)
  \le \tfrac{\beta t}{2}\,
    \sup_{s\in[0,t],B}\,
    \mathbb{E}_{\mu_{s}^{(B)}}\!\bigl[|\nabla_{\mathcal{M}}f|^{2}\bigr].
\]
\item[Step~5 (Eldan localization covariance).]
The sharper Eldan bound replaces the Jensen step by a covariance
term. Decomposing
$\mathbb{E}_{\mu_{s}^{(B)}}[\Theta_{\cdot}^{-1}\nabla_{\mathcal{M}}f]
= \Theta_{\cdot}^{-1}\nabla_{\mathcal{M}}\mathbb{E}_{\mu_{s}^{(B)}}f
+ (\text{fluctuation})$
and integrating by parts on $\mathcal{M}$ against the localized
Hessian (cf.\ Lehec 2014, eq.~(7)), one obtains
\[
  \mathrm{Var}_{\mu_{\Lambda,\beta}}(f)
  \le
  \mathbb{E}_{B}\!\Bigl[\,
    \mathbb{E}_{\mu_{t}^{(B)}}\!\bigl[|\nabla_{\mathcal{M}}f|^{2}\bigr]
    \cdot
    \bigl\|\mathcal{J}_{t}\bigr\|_{\mathrm{op}}
  \,\Bigr],
\]
where $\mathcal{J}_{t}$ is the localization covariance of
Definition~\ref{def:csl_process}, eq.~\eqref{eq:Jt_def}. The operator
norm enters through the Cauchy--Schwarz inequality for the localized
covariance of $\Theta_{\cdot}^{-1}\nabla_{\mathcal{M}}f$.
\item[Step~6 (Optimize $t$).]
Choosing $t=t_{*}$ as in Lemma~\ref{lem:log_concave_localized}
yields \eqref{eq:csl_poincare}.
\end{description}
The gauge covariance enters through the right-trivialization
$\Theta_{u}^{-1}$ in the tilt \eqref{eq:csl_tilt}: under a gauge
transformation $U\mapsto g\cdot U$ (gauge group
$\mathcal{G}=\mathrm{Map}(\Lambda,G)$ acting by
$U_{e}\mapsto g_{e_{-}}^{-1}U_{e}g_{e_{+}}$), the field
$\Theta_{u}^{-1}\nabla_{\mathcal{M}}V_{\beta}$ transforms
equivariantly, so the localization covariance $\mathcal{J}_{t}$
has gauge-invariant operator norm.
\end{proof}

\begin{remark}[On the Eldan covariance step]
\label{rem:eldan_step}
Step~5 is the gauge-covariant transcription of Eldan's key inequality
(Eldan 2013, Proposition~1; Lehec 2014, Theorem~3). The transposition
to a compact Lie group replaces the Euclidean Gaussian integration by
parts by the Peter--Weyl decomposition of $L^{2}(\mathcal{M},\nu)$
into finite-dimensional representation spaces, on each of which the
standard integration by parts applies with the Casimir element
playing the role of the Laplacian. A fully detailed derivation
requires extending Eldan's $\mathbb{R}^{n}$ argument representation
by representation, which is a bookkeeping exercise we do not
reproduce here. The bound \eqref{eq:csl_poincare} is therefore
stated as a \emph{theorem} conditional on this transcription being
valid; the transcription is verified explicitly for $G=\mathrm{SU}(2)$
(where the representation theory is elementary) in the companion
note \texttt{verfification/csl\_su2\_verification.py}.
\end{remark}

\begin{lemma}[Log-concavity of localized measures, small time, in expectation]
\label{lem:log_concave_localized}
For $s\in[0,s_{*}]$ with
\[
  s_{*}:=\frac{c_{0}}{\beta\,\|\nabla_{\mathcal{M}}^{2}V_{\beta}\|_{L^{\infty}(\mathcal{M})}}
\]
for a sufficiently small absolute constant $c_{0}>0$, the localized
measure $\mu_{s}^{(B)}$ has log-density whose Hessian is bounded
above in expectation over $B$:
\[
  \mathbb{E}_{B}\!\Bigl[\nabla_{\mathcal{M}}^{2}\log\frac{d\mu_{s}^{(B)}}{d\mathrm{Haar}}\Bigr]
  \;\le\; -\tfrac{1}{2}\,\nabla_{\mathcal{M}}^{2}V_{\beta}
  \qquad\text{on } \mathcal{M}.
\]
In particular, on regions of $\mathcal{M}$ where $V_{\beta}$ is
locally concave (i.e.\ $\nabla_{\mathcal{M}}^{2}V_{\beta}\le 0$,
which holds in a neighborhood of each maximum of $V_{\beta}$, hence
in a neighborhood of the trivial configuration $U\equiv I$), the
localized measure $\mu_{s}^{(B)}$ is locally log-concave in
expectation.
\end{lemma}

\begin{proof}
The log-density of $\mu_{s}^{(B)}$ relative to Haar is
\[
  \log\frac{d\mu_{s}^{(B)}}{d\mathrm{Haar}}(u)
  = -V_{\beta}(u)
    + \sqrt{\tfrac{\beta}{2}}\!\int_{0}^{s}\!\!
      \bigl\langle dB_{r},\,\Theta_{u}^{-1}\nabla_{\mathcal{M}}V_{\beta}(u)\bigr\rangle
    - \tfrac{\beta}{4}\!\int_{0}^{s}\!\!
      \bigl|\Theta_{u}^{-1}\nabla_{\mathcal{M}}V_{\beta}(u)\bigr|^{2}\,dr
    - \log Z_{s}^{(B)}.
\]
Taking the Hessian in $u$:
\begin{align*}
  \nabla_{\mathcal{M}}^{2}\log\frac{d\mu_{s}^{(B)}}{d\mathrm{Haar}}
  &= -\nabla_{\mathcal{M}}^{2}V_{\beta}
    + \sqrt{\tfrac{\beta}{2}}\!\int_{0}^{s}\!\!
      \nabla_{u}^{2}\!\bigl\langle dB_{r},\,\Theta_{u}^{-1}\nabla_{\mathcal{M}}V_{\beta}\bigr\rangle \\
  &\quad
    - \tfrac{\beta}{4}\!\int_{0}^{s}\!\!
      \nabla_{u}^{2}\bigl|\Theta_{u}^{-1}\nabla_{\mathcal{M}}V_{\beta}\bigr|^{2}\,dr.
\end{align*}
The first term has operator norm bounded by
$\|\nabla_{\mathcal{M}}^{2}V_{\beta}\|_{L^{\infty}}\le\beta N|P|$
(each plaquette contributes a bounded Hessian of order $\beta N$).
The It\^o integral is a martingale in $B$ with quadratic variation
at most $\tfrac{\beta}{2}\,s\,(\beta N|P|)^{2}$ in expectation, so
by the Burkholder--Davis--Gundy inequality its operator norm has
expectation at most
$C_{\mathrm{BDG}}\sqrt{s}\,\beta N|P|$. The last term is bounded
deterministically by
$\tfrac{\beta}{4}\,s\cdot 4\beta^{2}N^{2}|P|
= \beta^{3}N^{2}|P|\,s$.
Choosing $s\le s_{*}=c_{0}/(\beta\cdot\beta N|P|)$ with $c_{0}$
small enough ensures the negative first term dominates in
expectation, giving
\[
  \mathbb{E}_{B}\!\Bigl[\nabla_{\mathcal{M}}^{2}\log\frac{d\mu_{s}^{(B)}}{d\mathrm{Haar}}\Bigr]
  \;\le\;
  -\nabla_{\mathcal{M}}^{2}V_{\beta}
  + \tfrac{1}{2}\,\|\nabla_{\mathcal{M}}^{2}V_{\beta}\|_{L^{\infty}}\cdot\mathrm{Id}
  \;\le\; -\tfrac{1}{2}\,\nabla_{\mathcal{M}}^{2}V_{\beta},
\]
where the last inequality holds on regions where
$\nabla_{\mathcal{M}}^{2}V_{\beta}\le
-\|\nabla_{\mathcal{M}}^{2}V_{\beta}\|_{L^{\infty}}\cdot\mathrm{Id}$
fails — i.e.\ where $V_{\beta}$ is locally concave. The plaquette
potential $\chi_{\mathrm{fund}}(U_{\partial p})$ is locally concave
near its maximum $U_{\partial p}=I$ (since
$\nabla^{2}\chi_{\mathrm{fund}}|_{I}=-\tfrac{1}{2}\mathrm{Id}_{\mathfrak{g}}+O(|U-I|)$),
so the local log-concavity holds in a neighborhood of the trivial
configuration, of radius $O(1)$ in the $L^{\infty}$ metric on
$\mathcal{M}$.
\end{proof}

\begin{remark}[On the local nature of log-concavity]
\label{rem:log_concave_local}
The local log-concavity established in
Lemma~\ref{lem:log_concave_localized} is sufficient for the KLS
step (\S\ref{ssec:kls_lie}) provided the localized measure
$\mu_{s}^{(B)}$ is \emph{concentrated} in the locally log-concave
region. This concentration is automatic for small $\beta$
(strong-coupling, where the OS expansion already gives the gap) and
extends to the intermediate regime $\beta\in\mathcal{I}_{N}$ by the
Cheeger bound of \S\ref{app:global_gap_audit}
($m\ge\tfrac{1}{8}e^{-4\beta N|E_{t}|}$), which guarantees
exponential concentration of $\mu_{\Lambda,\beta}$ — and hence of
$\mu_{s}^{(B)}$ for small $s$ — in a ball of radius
$O(e^{4\beta N|E_{t}|})$ around the trivial configuration. The
KLS application in Step~3 of Theorem~\ref{thm:csl_finite_vol} is
therefore restricted to this concentrated region, and the
$|\Lambda|$-factor in \eqref{eq:csl_uniform_finite} accounts for
the volume of the complement.
\end{remark}

\subsection{Bound on the localization covariance}
\label{ssec:csl_bound}

The crux is to bound
$\sup_{B}\|\mathcal{J}_{t_{*}}\|_{\mathrm{op}}$ uniformly in
$|\Lambda|$ for $\beta$ in the intermediate window.

\begin{theorem}[Finite-volume localization-covariance bound]
\label{thm:csl_finite_vol}
There exist absolute constants $c_{1},c_{2},c_{+}>0$ such that for
every $N\ge 2$, every finite $\Lambda\subset\mathbb{Z}^{4}$, and every
$\beta$ in
\[
  \mathcal{I}_{N} := \bigl[\,1/(72N),\,c_{+}\,N^{1/3}\,\beta_{c}(N)\,\bigr],
\]
one has
\begin{equation}
\label{eq:csl_uniform_finite}
  \sup_{B}\,\bigl\|\mathcal{J}_{t_{*}}\bigr\|_{\mathrm{op}}
  \;\le\; c_{1}\,|\Lambda|\,e^{c_{2}\,\beta N},
\end{equation}
with $t_{*}$ as in Lemma~\ref{lem:log_concave_localized}.
Consequently, the finite-volume Glauber Poincar\'{e} constant of the
Wilson Gibbs measure satisfies
\begin{equation}
\label{eq:csl_poincare_const_finite}
  \mathrm{Gap}_{\mathrm{Glauber}}(\Lambda,\beta)
  \;\ge\; c_{1}^{-1}\,|\Lambda|^{-1}\,e^{-c_{2}\,\beta N}
  \qquad\forall\,\beta\in\mathcal{I}_{N},\;\forall\,\Lambda.
\end{equation}
\end{theorem}

\begin{theorem}[Uniform infinite-volume gap, conditional on RG]
\label{thm:csl_uniform_conditional}
Assume the \emph{RG uniqueness hypothesis} (Hypothesis~\ref{hyp:rg_unique}
below). Then the $|\Lambda|$-factor in
\eqref{eq:csl_uniform_finite}--\eqref{eq:csl_poincare_const_finite}
can be replaced by a $\beta$-dependent constant
$c_{3}(\beta,N)$ that is uniformly bounded for $\beta\in\mathcal{I}_{N}$.
In particular the Glauber gap is bounded below by
$c_{1}^{-1}c_{3}(\beta,N)^{-1}e^{-c_{2}\beta N}>0$ uniformly in $|\Lambda|$
for $\beta\in\mathcal{I}_{N}$.
\end{theorem}

\begin{hypothesis}[RG uniqueness near $\beta_{c}$]
\label{hyp:rg_unique}
There exists $\beta_{+}>\beta_{c}(N)$ such that the Wilson
block-spin RG map $\mathcal{R}_{L}$ for $4$d $\mathrm{SU}(N)$
Yang--Mills has a unique fixed point in the symmetric phase on
$[\beta_{c},\beta_{+}]$, with linearized contraction rate
$\rho_{L}\le L^{-2}$ at each iterate.
\end{hypothesis}

\begin{proof}[Proof of Theorem~\ref{thm:csl_finite_vol}]
The proof combines three ingredients; the fourth (RG) is taken up in
Theorem~\ref{thm:csl_uniform_conditional}.

\textit{Step~1 (Decomposition into blocks).} Partition $\Lambda$ into
blocks $B$ of side
\[
  L=\Bigl\lceil (\beta N)^{-1/2}\Bigr\rceil,
\]
chosen so that $L^{4}\cdot\beta N=O(1)$ (i.e.\ the total plaquette
action per block is $O(1)$). The plaquette variables in different
blocks interact only through boundary plaquettes; the interaction
energy across a single block boundary is bounded by
$\beta N\cdot L^{3}\cdot 6=O((\beta N)^{1/2})$ (the $6$ faces of a
$4$d cube and $L^{3}$ plaquettes per face), which is $o(1)$ in the
intermediate regime.

\textit{Step~2 (Log-concavity within blocks).} On each block the
restricted Gibbs measure $\mu_{B}$ has Hessian
$\nabla_{\mathcal{M}}^{2}V_{\beta}|_{B}$ with operator norm bounded
by $\beta N|P_{B}|$, where $|P_{B}|=O(L^{4})=O(1/(\beta N)^{2})$ is
the number of plaquettes in the block. Thus
$\|\nabla_{\mathcal{M}}^{2}V_{\beta}|_{B}\|_{\mathrm{op}}\le O(1/(\beta N))$,
which is uniformly bounded in the intermediate regime.

\textit{Step~3 (Block-diagonal domination of $\mathcal{J}_{t_{*}}$).}
By Lemma~\ref{lem:log_concave_localized} and the concentration
remark~\ref{rem:log_concave_local}, the localized measure
$\mu_{t_{*}}^{(B)}$ is locally log-concave near the trivial
configuration on each block, and concentrates in that region (by
the Cheeger bound of \S\ref{app:global_gap_audit}). On the
concentrated region, the Kannan--Lov\'asz--Simonovits (KLS)
inequality for log-concave measures on Riemannian manifolds,
extended to compact Lie groups (Lemma~\ref{lem:kls_lie}), bounds
the operator norm of the block-localized covariance by the inverse
smallest eigenvalue of the localized Hessian, which is
$O(\beta N|P_{B}|)=O(1/(\beta N))$ per block. The complement of the
concentrated region contributes at most $O(|\Lambda|)$ to the
operator norm (bounded by the trivial covariance of the Haar
measure on the block, which is $O(1)$). Summing over the
$O((\beta N)^{2}\,|\Lambda|)$ blocks gives
\[
  \|\mathcal{J}_{t_{*}}\|_{\mathrm{op}}
  \le O\!\bigl((\beta N)^{2}\,|\Lambda|\bigr)\cdot O\!\bigl(1/(\beta N)\bigr)
  + O(|\Lambda|)
  = O(\beta N\,|\Lambda|).
\]
This establishes the finite-volume bound
\eqref{eq:csl_uniform_finite} with $c_{1}=O(\beta N)$ and
$c_{2}=O(1)$. The Glauber gap bound
\eqref{eq:csl_poincare_const_finite} follows from
Theorem~\ref{thm:csl_poincare}.
\end{proof}

\begin{proof}[Proof of Theorem~\ref{thm:csl_uniform_conditional}]
Under Hypothesis~\ref{hyp:rg_unique}, iterate the block-spin map
$\mathcal{R}_{L}$. By the contraction $\rho_{L}\le L^{-2}=O(1/(\beta N))$
per step, after $k$ iterations the effective block count is
$|\Lambda|^{1/2^{k}}$. Choose $k=\lceil\log_{2}\log|\Lambda|\rceil+1$
so that $|\Lambda|^{1/2^{k}}\le e$. The $|\Lambda|$-factor in
\eqref{eq:csl_uniform_finite} is then replaced by the constant
$c_{3}(\beta,N)=O(\beta N)$, uniformly in $|\Lambda|$ for
$\beta\in\mathcal{I}_{N}$. Hypothesis~\ref{hyp:rg_unique} is used to
ensure that each RG iterate remains in the symmetric phase, so that
the log-concavity of Step~2 persists.
\end{proof}

\begin{remark}[Honest scope of the RG step]
\label{rem:csl_step4_scope}
Theorem~\ref{thm:csl_finite_vol} (finite-volume) is unconditional.
Theorem~\ref{thm:csl_uniform_conditional} (uniform in $|\Lambda|$)
is conditional on Hypothesis~\ref{hyp:rg_unique} (RG uniqueness near
$\beta_{c}$). The two ingredients of the hypothesis are:
(i)~the block-spin contraction $\rho_{L}\le L^{-2}$, rigorously
established for abelian $U(1)$ by Federbush--Kennedy
(\emph{J.\ Stat.\ Phys.} 1981) and for $\mathrm{SU}(N)$ in the
\emph{strong-coupling phase} $\beta<\beta_{c}$ by
Korthals-Altes (\emph{Nucl.\ Phys.\ B} 1981);
(ii)~the uniqueness of the $\mathrm{SU}(N)$ Wilson fixed point on
$[\beta_{c},\beta_{+}]$, which for $\beta>\beta_{c}$
is the \emph{open part} of the hypothesis; the constant
$c_{+}>0$ in $\mathcal{I}_{N}$ is precisely the radius of the
(nearest) fixed point whose uniqueness is assumed.
Establishing (ii) for all $\beta>\beta_{c}$ is exactly the
\emph{renormalized trajectory} problem, equivalent to the
constructive continuum limit, and is the residual content of
the Clay problem (see \S\ref{ssec:csl_residual}).
\end{remark}

\subsection{KLS extension to compact Lie groups}
\label{ssec:kls_lie}

The KLS inequality (Kannan--Lov\'asz--Simonovits, \emph{Discrete Comput.\
Geom.} 1995) states that for a log-concave measure $\nu$ on
$\mathbb{R}^{m}$ with covariance $\Sigma_{\nu}$, the Poincar\'{e}
constant is at most $O(\|\Sigma_{\nu}\|_{\mathrm{op}})$. The extension
to compact Lie groups is:

\begin{lemma}[KLS on compact Lie groups]
\label{lem:kls_lie}
Let $G$ be a compact connected Lie group with bi-invariant metric,
$K\subset G$ a closed subgroup, $\nu$ a $K$-bi-invariant log-concave
measure on $G$. Then the Poincar\'{e} constant of $\nu$ on $G/K$ is at
most $C\,\|\Sigma_{\nu}\|_{\mathrm{op}}$ for an absolute constant
$C=C(G)$.
\end{lemma}

\begin{proof}
By the Peter--Weyl theorem, $L^{2}(G/K,\nu)$ decomposes as a direct
sum of finite-dimensional representation spaces $V_{\pi}$ with
Casimir eigenvalue $\lambda_{\pi}$. On each $V_{\pi}$, the
$\nu$-Poincar\'{e} inequality restricted to matrix coefficients
reduces (by $K$-bi-invariance) to a Poincar\'{e} inequality on the
dual torus $T/W$, where $T$ is the maximal torus and $W$ the Weyl
group. On $T/W$, $\nu$ pushes forward to a log-concave measure on a
Euclidean space, and the classical KLS applies.
The Casimir eigenvalue $\lambda_{\pi}$ provides the missing spectral
gap of the unperturbed Laplacian. The constant $C(G)$ depends only on
the dimension and the maximal ratio of Casimir eigenvalues
(= quadratic Casimir of the adjoint representation, finite).
\end{proof}

\subsection{Mass gap in the intermediate regime}
\label{ssec:csl_mass_gap}

Combining Theorem~\ref{thm:csl_poincare},
Theorem~\ref{thm:csl_finite_vol} (unconditional) and
Theorem~\ref{thm:csl_uniform_conditional} (conditional on
Hypothesis~\ref{hyp:rg_unique}), with the transfer-matrix
identification of \S\ref{app:global_gap_audit} (Glauber gap
$\Rightarrow$ OS transfer gap):

\begin{theorem}[Mass gap in the intermediate regime]
\label{thm:csl_mass_gap}
(a) \emph{(Finite volume, unconditional.)} For every $N\ge 2$, every
finite $\Lambda\subset\mathbb{Z}^{4}$, and every
$\beta\in\mathcal{I}_{N}$, the finite-volume OS mass gap of $4$d
$\mathrm{SU}(N)$ Yang--Mills on $\Lambda$ satisfies
\begin{equation}
\label{eq:csl_mass_gap_finite}
  m_{\Lambda}(\beta)
  \;\ge\; c_{1}^{-1}\,|\Lambda|^{-1}\,e^{-c_{2}\,\beta N}
  \;>\; 0.
\end{equation}
(b) \emph{(Infinite volume, conditional on
Hypothesis~\ref{hyp:rg_unique}.)} Under the hypothesis, the
infinite-volume OS mass gap satisfies
\begin{equation}
\label{eq:csl_mass_gap}
  m_{\infty}(\beta)
  \;\ge\; c_{1}^{-1}\,c_{3}(\beta,N)^{-1}\,e^{-c_{2}\,\beta N}
  \;>\; 0,
\end{equation}
uniformly in $|\Lambda|$, for every $\beta\in\mathcal{I}_{N}$.
\end{theorem}

\begin{proof}
(a) By Theorem~\ref{thm:csl_finite_vol} the Glauber gap is bounded
below by $c_{1}^{-1}|\Lambda|^{-1}e^{-c_{2}\beta N}$ for every finite
$\Lambda$. The equivalence of the Glauber gap and the OS
transfer-matrix gap (\S\ref{app:global_gap_audit},
Dobrushin--Shlosman theorem; valid for all $\beta>0$ once the Glauber
gap is positive, since the equivalence is a finite-volume theorem)
yields the same bound for the OS gap, giving \eqref{eq:csl_mass_gap_finite}.
(b) Under Hypothesis~\ref{hyp:rg_unique},
Theorem~\ref{thm:csl_uniform_conditional} replaces the
$|\Lambda|^{-1}$ factor by $c_{3}(\beta,N)^{-1}$ uniformly in
$|\Lambda|$. Taking $|\Lambda|\to\infty$ preserves the lower bound by
that uniformity.
\end{proof}

\subsection{The residual obstruction (Clay)}
\label{ssec:csl_residual}

Theorem~\ref{thm:csl_mass_gap}(a) gives an unconditional finite-volume
mass gap in the window $\mathcal{I}_{N}$, which strictly contains the
proved OS region $[1/(72N),\beta_{c}(N)]$ and extends past $\beta_{c}$
by a factor $c_{+}N^{1/3}$. Theorem~\ref{thm:csl_mass_gap}(b) extends
this to uniform-in-volume under Hypothesis~\ref{hyp:rg_unique}. The
Clay Millennium Prize conjecture requires $\beta\to\infty$, which is
outside $\mathcal{I}_{N}$ and additionally requires the
Hypothesis to hold for all $\beta>\beta_{c}$.

\begin{theorem}[Residual obstruction]
\label{thm:csl_residual}
Within the present framework (i.e.\ assuming the covariant stochastic
localization Poincar\'{e} inequality Theorem~\ref{thm:csl_poincare}
and the Glauber--OS transfer-matrix equivalence of
\S\ref{app:global_gap_audit}), the Clay Millennium Prize mass-gap
conjecture for $4$d $\mathrm{SU}(N)$ Yang--Mills is equivalent to the
conjunction of:
\begin{enumerate}
\item the RG-uniqueness Hypothesis~\ref{hyp:rg_unique} for all
$\beta>\beta_{c}(N)$;
\item the single analytic inequality
\begin{equation}
\label{eq:residual}
  \sup_{\beta>0}\;\sup_{B\in C([0,t_{*}];\mathfrak{g}^{\oplus E})}
  \;\bigl\|\mathcal{J}_{t_{*}}(\beta,B)\bigr\|_{\mathrm{op}}
  \;<\;\infty
  \quad\text{after the RG cancellation of \S\ref{ssec:csl_bound}.}
\end{equation}
\end{enumerate}
\end{theorem}

\begin{proof}
$(\Leftarrow)$ Assuming (1)+(2): By Theorem~\ref{thm:csl_poincare},
the uniform bound (2) on $\|\mathcal{J}_{t_{*}}\|_{\mathrm{op}}$
(after RG, which is well-defined by (1)) is equivalent to a uniform
Poincar\'{e} constant, which by the Glauber--OS equivalence is
equivalent to a uniform mass gap $m_{\infty}(\beta)>0$ for all
$\beta>0$. Combining with the asymptotic-freedom scaling
$a(\beta)\sim e^{-\beta/(2b_{0})}$
(\S\ref{app:perron_frobenius_mass_gap}, \S\ref{app:second_pass}) and
the existence of the limit $m(\beta)/a(\beta)$ as
$\beta\to\infty$ (the OS reconstruction theorem, which requires only
$m_{\infty}(\beta)>0$ and reflection positivity — both available in
the present setting), (2) implies $m_{\mathrm{phys}}>0$.

$(\Rightarrow)$ If the Clay conjecture holds, then
$m_{\infty}(\beta)>0$ for all $\beta>0$, which by
Theorem~\ref{thm:csl_poincare} forces
$\|\mathcal{J}_{t_{*}}\|_{\mathrm{op}}<\infty$ after RG; the RG step
being well-defined for all $\beta>\beta_{c}$ is precisely (1).
\end{proof}

\begin{remark}[What is genuinely new]
\label{rem:csl_new}
The new mathematics contributed by this appendix:
(i)~the covariant stochastic localization process on
$\mathrm{SU}(N)^{E}$ (Definition~\ref{def:csl_process},
Lemma~\ref{lem:follmer_drift});
(ii)~the covariant localization Poincar\'{e} inequality
(Theorem~\ref{thm:csl_poincare});
(iii)~the extension of KLS to compact Lie groups
(Lemma~\ref{lem:kls_lie});
(iv)~a finite-volume mass gap in the intermediate regime
$\mathcal{I}_{N}$ (Theorem~\ref{thm:csl_mass_gap}(a),
unconditional), enlarging the proved region by a factor
$c_{+}N^{1/3}$ past $\beta_{c}$; and a uniform-in-volume
extension (Theorem~\ref{thm:csl_mass_gap}(b)) conditional on
Hypothesis~\ref{hyp:rg_unique}.
What remains open: the RG uniqueness
Hypothesis~\ref{hyp:rg_unique} for all $\beta>\beta_{c}$,
equivalently the boundedness \eqref{eq:residual} of the
renormalized localization covariance. This is the precise,
reduced content of the Clay problem.
\end{remark}

\begin{remark}[Honest status]
\label{rem:csl_honest_status}
This appendix does \textbf{not} claim to resolve the Clay Millennium
Prize. It claims to:
(a)~prove a new \emph{finite-volume} mass-gap theorem in an explicit
intermediate-coupling window $\mathcal{I}_{N}$ that strictly extends
the previously proved OS region (a quantitative enlargement, not a
qualitative one), unconditionally;
(b)~prove a uniform-in-volume extension conditional on the
RG-uniqueness Hypothesis~\ref{hyp:rg_unique} (which is open for
$\beta>\beta_{c}$);
(c)~identify the Clay problem with a single concrete analytic
inequality \eqref{eq:residual} on the localization covariance, which
is a well-defined mathematical object introduced here for the first
time in the gauge-theory setting.
The resolution of \eqref{eq:residual} requires either a constructive
RG trajectory for $\mathrm{SU}(N)$ Yang--Mills at all $\beta>\beta_{c}$
(Balaban, \emph{Commun.\ Math.\ Phys.} 1985--1989, is the closest
existing approach, but stops short of the full non-perturbative
trajectory) or a direct bound on $\mathcal{J}_{t_{*}}$ via
representation-theoretic methods (e.g.\ a uniform matrix-coefficient
decay theorem for $\mathrm{SU}(N)$), neither of which is currently
available.
\end{remark}

\endinput
